\documentclass[UTF8]{ctexart}

\usepackage{amsmath, amssymb}
\usepackage{times}
\usepackage[a4paper, landscape, left=1.2cm, right=1.2cm, top=1.2cm, bottom=1.2cm, headheight=14pt, headsep=6pt]{geometry}
\usepackage{fancyhdr}
\usepackage{xcolor}
\usepackage{tabularx}
\usepackage{booktabs}
\usepackage{tikz}
\usetikzlibrary{shapes.geometric, arrows.meta, positioning, calc}
\usepackage[hidelinks]{hyperref}

% =====================================================================================
% 颜色配置与页面定义
% =====================================================================================
\definecolor{themeblue}{RGB}{31, 78, 121}
\definecolor{themelightblue}{RGB}{221, 235, 247}
\definecolor{alertred}{RGB}{192, 0, 0}
\definecolor{fpgagreen}{RGB}{38, 128, 0}
\definecolor{gategold}{RGB}{255, 242, 204}
\definecolor{apppurple}{RGB}{230, 224, 236}
\definecolor{appborder}{RGB}{112, 48, 160}
\definecolor{algogreen}{RGB}{226, 240, 217}
\definecolor{algoborder}{RGB}{84, 130, 53}

% =====================================================================================
% 超链接徽章宏定义 (双向绑定图表与表格)
% =====================================================================================
\newcommand{\bdg}[1]{\hypertarget{flow#1}{\hyperlink{tab#1}{\tikz[baseline=-0.6ex]\node[circle, fill=alertred, text=white, inner sep=1pt, minimum size=3.5mm, font=\sffamily\bfseries\tiny] {#1};}}}
\newcommand{\tabbdg}[1]{\hypertarget{tab#1}{\hyperlink{flow#1}{\tikz[baseline=-0.6ex]\node[circle, fill=alertred, text=white, inner sep=1pt, minimum size=3.5mm, font=\sffamily\bfseries\tiny] {#1};}}}

\pagestyle{fancy}
\fancyhf{}
\fancyhead[L]{\textbf{嵌入式 Linux 跨部门协同研发与调试工程流程规范}}
\fancyhead[R]{\textbf{Ver 1.1}}
\fancyfoot[C]{\thepage}

\begin{document}

% =====================================================================================
% PAGE 1: 流程图 (保留原版 4行 S-Curve 拓扑布局 + 3.5层算法节点)
% =====================================================================================
\begin{center}
    {\Large \textbf{\color{themeblue} 嵌入式 Linux 软硬件全周期协同研发与调试流程}} \\[0.2em]
    {\scriptsize \textbf{参与团队：} 硬件(EE)组、FPGA组、ARM组、算法(Algo)组、测试(QA)组 \\ 
    \textit{*提示：点击图表或表格中的 \bdg{1} 红色标号，可实现双向超链接跳转。}}
\end{center}

\vspace{-0.1cm}

\begin{figure}[!ht]
\centering
\resizebox{0.98\linewidth}{!}{%
\begin{tikzpicture}[
    >=Stealth,
    every text node part/.style={align=center},
    startstop/.style={rectangle, rounded corners=3pt, minimum width=2.4cm, minimum height=0.85cm, fill=themeblue!90, text=white, font=\bfseries\scriptsize},
    process/.style={rectangle, rounded corners=2pt, minimum width=2.4cm, minimum height=0.85cm, fill=themelightblue, draw=themeblue, line width=0.8pt, font=\scriptsize},
    gate/.style={diamond, aspect=1.9, minimum width=2.5cm, minimum height=0.85cm, fill=gategold, draw=orange!90!black, line width=0.8pt, font=\scriptsize\bfseries},
    appnode/.style={rectangle, rounded corners=2pt, minimum width=2.4cm, minimum height=0.85cm, fill=apppurple, draw=appborder, line width=0.8pt, font=\scriptsize},
    algonode/.style={rectangle, rounded corners=2pt, minimum width=2.4cm, minimum height=0.85cm, fill=algogreen, draw=algoborder, line width=0.8pt, font=\scriptsize},
    feedback/.style={rectangle, rounded corners=2pt, fill=red!10, draw=alertred, line width=0.8pt, font=\tiny\bfseries}
]

    % ----------------- ROW 1: 硬件设计与联合评估 (向右 ->) -----------------
    \node (start)       at (0, 0)      [startstop] {\bdg{1} 需求与任务书\\[-1pt]\tiny 软硬件同步输入};
    \node (ee_audit)    at (4.5, 0)    [gate]      {\bdg{2} 硬件(EE)选型\\[-1pt]\tiny 驱动/Pinmux审查};
    \node (fpga_audit)  at (9.0, 0)    [gate, dashed]      {\bdg{3} [可选] FPGA/MCU协议\\[-1pt]\tiny 软/硬/算法联合审查};
    \node (sch_pcb)     at (13.5, 0)   [process]   {\bdg{4} 原理图与 Layout\\绘制及联合评审};

    % Row 1 专用报错节点
    \node (ee_reject)   at (2.25, 1.6)  [feedback] {驳回选型\\EE组重新选型};
    \node (fpga_reject) at (6.75, 1.6)  [feedback] {驳回协议\\重新协商制订};

    % ----------------- ROW 2: 底层与中间件 Bring-up (向左 <-) -----------------
    \node (pcba_fab)    at (13.5, -2.8) [process]   {\bdg{5} PCBA 打样贴片\\与阻抗初检};
    \node (hw_bringup)  at (9.0, -2.8)  [gate]      {\bdg{6} 裸板硬件上电\\PMIC 时序测量};
    \node (bsp_os)      at (4.5, -2.8)  [process]   {\bdg{7} SPL/Kernel 启动\\(ARM组)};
    \node (driver_fpga) at (0, -2.8)    [gate]      {\bdg{8} 外设/FPGA联调\\(ARM组)};

    % 硬件 ECN 报错节点 (外置于最右侧专属通道)
    \node (ecn_fix)     at (16.2, -1.4) [feedback] {硬件 ECN 变更\\改板重新打样};

    % ----------------- ROW 3: 上层业务开发与黑盒测试 (向右 ->) -----------------
    \node (app_ui)      at (0, -5.6)    [appnode]   {\bdg{9.1} 应用与 UI 开发\\(ARM组)};
    \node (board_screen)at (4.5, -5.6)  [process]   {\bdg{10} 板筛固件与大纲\\(产测准备)};
    \node (qa_test)     at (9.0, -5.6)  [gate]      {\bdg{11} 黑盒集成测试\\(QA测试组)};
    \node (ess_test)    at (13.5, -5.6) [gate]      {\bdg{12} 环境应力筛选\\高低温/振动压测};

    % ----------------- ROW 3.5: 算法集成专属层 -----------------
    \node (algo_int)    at (0, -7.5)    [algonode]  {\bdg{9.2} 算法集成\\(算法组)};

    % 内部逻辑报错置于专属夹层 (Y=-4.2 与 Y=-8.6)
    \node (bug_fix)     at (4.5, -4.2)  [feedback] {测试Bug修复\\[-1pt]\tiny 版本迭代优化};
    \node (ess_fail)    at (6.75, -8.6) [feedback] {温漂/底层崩溃\\触发硬件/ARM回溯};

    % ----------------- ROW 4: 资料归档 -----------------
    \node (archiving)   at (13.5, -9.8) [startstop] {\bdg{13} 资料归档转产\\(SDK/BOM/SOP)};

    % ================= 连线逻辑 (Strictly Manhattan Orthogonal Routing) =================
    
    % --- Main Flow: Row 1 (L->R) ---
    \draw [->, line width=0.8pt] (start.east) -- node[above=1pt, font=\tiny\bfseries] {并行分解} (ee_audit.west);
    \draw [->, line width=0.8pt] (ee_audit.east) -- node[above=1pt, font=\tiny\bfseries] {通过} (fpga_audit.west);
    \draw [->, line width=0.8pt] (fpga_audit.east) -- node[above=1pt, font=\tiny\bfseries] {无FPGA或 \\ 评审通过} (sch_pcb.west);

    % Transition Row 1 -> Row 2
    \draw [->, line width=0.8pt] (sch_pcb.south) -- node[right=1pt, font=\tiny\bfseries] {评审通过} (pcba_fab.north);

    % --- Main Flow: Row 2 (R->L) ---
    \draw [->, line width=0.8pt] (pcba_fab.west) -- node[above=1pt, font=\tiny\bfseries] {交接硬件} (hw_bringup.east);
    \draw [->, line width=0.8pt] (hw_bringup.west) -- node[above=1pt, font=\tiny\bfseries] {时序正常} (bsp_os.east);
    \draw [->, line width=0.8pt] (bsp_os.west) -- node[above=1pt, font=\tiny\bfseries] {系统点亮} (driver_fpga.east);

    % Transition Row 2 -> Row 3 (API就绪与反向需求修改)
    \draw [->, line width=0.8pt] ([xshift=-5pt]app_ui.north) -- node[left=1pt, font=\tiny\bfseries] {按需修改 API \\ 或FPGA/MCU协议} ([xshift=-5pt]driver_fpga.south);
    \draw [->, line width=0.8pt] ([xshift=5pt]driver_fpga.south) -- node[right=1pt, font=\tiny\bfseries] {API \\ 就绪} ([xshift=5pt]app_ui.north);

    % --- Main Flow: Row 3 (L->R) ---
    \draw [->, line width=0.8pt] (app_ui.east) -- node[above=1pt, font=\tiny\bfseries] {业务调通} (board_screen.west);
    \draw [->, line width=0.8pt] (board_screen.east) -- node[above=1pt, font=\tiny\bfseries] {大纲定版} (qa_test.west);
    \draw [->, line width=0.8pt] (qa_test.east) -- node[above=1pt, font=\tiny\bfseries] {用例通过} (ess_test.west);

    % Transition Row 3.5 (算法集成)
    \draw [<->, line width=0.8pt] (app_ui.south) -- node[right=1pt, font=\tiny\bfseries] {集成 .h/.so} (algo_int.north);

    % Transition Row 3 -> Row 4
    \draw [->, line width=0.8pt] ([xshift=5pt]ess_test.south) -- node[right=1pt, font=\tiny\bfseries] {压测通过} ([xshift=5pt]archiving.north);

    % ================= 异常驳回路径 (Red Orthogonal Lines) =================

    % FB1: EE Reject
    \draw [->, line width=0.8pt, color=alertred] ([xshift=-5pt]ee_audit.north) |- (ee_reject.east);
    \draw [->, line width=0.8pt, color=alertred] (ee_reject.west) -| (start.north);

    % FB2: FPGA/MCU Reject
    \draw [->, line width=0.8pt, color=alertred] (fpga_audit.north) |- (fpga_reject.east);
    \draw [->, line width=0.8pt, color=alertred] (fpga_reject.west) -| ([xshift=5pt]ee_audit.north);

    % FB3: Hardware ECN
    \draw [->, line width=0.8pt, color=alertred] ([xshift=5pt]hw_bringup.south) -- +(0, -0.6) -| (ecn_fix.south);
    \draw [->, line width=0.8pt, color=alertred] (ecn_fix.north) |- (sch_pcb.east);

    % FB4: QA Bug Fix
    \draw [->, line width=0.8pt, color=alertred] ([xshift=-5pt]qa_test.north) |- (bug_fix.east);
    \draw [->, line width=0.8pt, color=alertred] (bug_fix.west) -| ([xshift=25pt]app_ui.north);

    % FB5: ESS Environmental Failure (走底部的 Y=-8.6 水平通道)
    \draw [->, line width=0.8pt, color=alertred] ([xshift=-5pt]ess_test.south) |- (ess_fail.east);
    \draw [->, line width=0.8pt, color=alertred] (ess_fail.west) -- (-3.6, -8.6) |- ([yshift=0pt]driver_fpga.west);

\end{tikzpicture}
}
\end{figure}

\newpage
% =====================================================================================
% PAGE 2: 质量控制表 (表1：输入依赖与门禁机制 | 表2：工作内容、产出与风险)
% =====================================================================================

\begin{center}
    {\large \textbf{\color{themeblue} 表 1：项目依赖关系}}
\end{center}
\vspace{-0.2cm}

\begin{table}[!ht]
\centering
\scriptsize
\renewcommand{\arraystretch}{1.25}
\begin{tabularx}{\linewidth}{@{} c l X X @{}}
\toprule
\textbf{序号} & \textbf{所处节点} & \textbf{前置依赖与触发输入项 (Inputs)} & \textbf{质量门禁与驳回机制 (Gates \& Goto)} \\ 
\midrule
\tabbdg{1} & 需求与任务书 & PM下发《产品设计定义》、《软件设计需求说明书》等。 & 若硬件选项无法承载软件业务设计，一票否决驳回重新评估立项。 \\
\tabbdg{2} & 硬件(EE)选型 & 依赖基础需求，EE 组输出初步芯片方案与引脚分配。 & 芯片须具备稳定主线驱动及原厂支持；否则门禁不通过驳回 EE 重选。 \\
\tabbdg{3} & FPGA/MCU协议 & 依赖传感器/外设接口等关键硬件选型初定 (可选节点)。 & 确认软/硬/算法通讯逻辑完备；若时序相位或控制流有缺陷则驳回重配。 \\
\tabbdg{4} & 原理图 Layout& 依赖 EE 组向 ARM 组输出提供完整的 PDF 原理图。 & 核心引脚严禁冲突，功能接口需要满足要求；否则阻断发板，驳回改图。 \\
\tabbdg{5} & PCBA 打样贴片 & 依赖 Layout 定版工程文件 (Gerber) 及 BOM 清单表。由 EE 组同事保障。 & 基本的空板阻抗排查测试通过，且元件无虚焊、无大面积连锡即可交接。 \\
\tabbdg{6} & 裸板硬件上电 & 依赖 SM 产线或手工焊接的首件 PCBA 实物产出交付。由 EE 组同事保障。 & 供电纹波合格；若时序严重错乱或存在不可修复的短路，触发 ECN 驳回改板。 \\
\tabbdg{7} & Kernel启动 & 依赖硬件核心电源树时序正常，且主控芯片及主频晶振成功起振。 & DDR 压测零报错；内核启动无 Oops/Panic；必须保障基本的串口控制台交互正常。 \\
\tabbdg{8} & 外设/FPGA联调& 依赖底层系统基础点亮，文件系统 (RootFS) 运行环境搭建完毕。 & 图像无撕裂；能向下稳定提供控制 API；否则要求底层驱动组重新修正时序或协议。 \\
\tabbdg{9.1} & 应用与UI开发 & 依赖 ARM 底层驱动提供的高层控制 API 及算法组交付的动态库。 & 业务主干逻辑闭环；若联调时接口抛出异常或吞吐率不达标，立即驳回要求修正底层。 \\
\tabbdg{9.2} & 算法集成 & 依赖算法组提供的完整 `.h` 头文件及适用于目标平台的 `.so` 库。 & 算法耗时与实际内存占用必须符合项目设计达标；否则算法组须主动配合联调优化。 \\
\tabbdg{10} & 板筛固件大纲 & 依赖基础业务逻辑全线打通，各项外设驱动功能均可控可达。 & 核心产测项覆盖率强制达 100\%；必须具备供量产车间使用的自动化一键 PASS/FAIL 功能。 \\
\tabbdg{11} & 黑盒集成测试 & 依赖定版板筛固件冻结受控，及 QA 审查确认通过后的《测试大纲》。 & P0/P1 用例全部 PASS；若发生死机崩溃或严重交互错误，立即开单驳回要求研发组修复 Bug。 \\
\tabbdg{12} & 环境应力(ESS)& 依赖 QA 组对于当前软硬件版本的集成功能测试全量验证通过结论。 & 极限温区运行无亚稳态现象；若发生硬死机灾难，无条件触发软硬件底层深度联合回溯调查。 \\
\tabbdg{13} & 资料归档转产 & 依赖最终的可靠性 ESS 环境极限压测分析报告结论被判定为合格。 & 软硬件版本号实现唯一性强绑定；各部门相关节点评审文档签署齐备；正式获批流转进入量产。 \\
\bottomrule
\end{tabularx}
\end{table}

\newpage

\vspace{0.4cm}

\begin{center}
    {\large \textbf{\color{themeblue} 表 2：具体工作内容、标准内容与潜在风险}}
\end{center}
\vspace{-0.2cm}

\begin{table}[!ht]
\centering
\scriptsize
\renewcommand{\arraystretch}{1.2}
\begin{tabularx}{\linewidth}{@{} c l X >{\raggedright\arraybackslash}p{3.5cm} X @{}}
\toprule
\textbf{序号} & \textbf{所处节点} & \textbf{具体工作内容 (Work Content)} & \textbf{输出交付件} & \textbf{潜在风险 (Potential Risks)} \\ 
\midrule
\tabbdg{1} & 需求与任务书 & 软、硬件同步并行分解需求，划定架构边界。明确具体架构逻辑，实施阶段\textbf{不承接}大范围变更。 & 《软件设计架构与功能》。 & 需求变更频繁导致架构被推翻，严重影响整体研制进度。 \\
\tabbdg{2} & 硬件(EE)选型 & ARM组同步评估芯片驱动情况。对内先作基础约束要求，后续视推进情况逐步新增细节审查项。 & 《软硬件选型评估表》。 & 核心元器件缺货，或缺少原厂/开源社区的技术与驱动支持。 \\
\tabbdg{3} & FPGA/MCU协议 & 联合 FPGA/算法组 制定控制逻辑与通信协议。仅在带自研前置处理的项目中联合输出。 & 《FPGA接口传输协议书》\cite{ref:protocol}。 & \textbf{算法组认为电机控制逻辑需补充，导致后端向前端提修改需求，造成返工。} \\
\tabbdg{4} & 原理图 Layout & 跨部门审查引脚冲突与电源域等配置。需引入内部 Checklist 机制逐一勾兑引脚，避免低级错漏。 & 《硬件设计评审表》。 & 高速信号阻抗不匹配或电源平面割裂，导致后期联调通信丢包。 \\
\tabbdg{5} & PCBA 贴片 & 贴片加工及首件目检。记录初级加工缺陷，防范批量转产阶段的 SMT 生产规模事故。 & 由 EE 组同事输出相关报告。 & SMT 贴片发生不可见 BGA 虚焊或大面积连锡，导致排查困难。 \\
\tabbdg{6} & 裸板硬件上电 & 测量 PMIC 上电时序及复位信号。持续积累裸板调试故障案例，为后续研制提供经验。 & 由 EE 组同事输出相关报告。 & 电源时序错乱或关键网络短路，导致一上电即烧毁核心 SoC 芯片。 \\
\tabbdg{7} & Kernel启动 & 移植 Bootloader，适配设备树(DTS)并挂载根文件系统。附带 Memory Map 与启动耗时时序分析。 & 系统镜像、源码仓库和《系统设计指南》\cite{ref:sysguide}。 & DDR 标定不稳定，导致在后续的高低温等恶劣环境下出现随机死机。 \\
\tabbdg{8} & 外设/FPGA联调 & 调试 DRM/I2C 等外设，按需响应算法底层修改。接口定义必须随代码同步更新，严禁口头联调。 & 系统镜像、源码仓库和《系统设计指南》\cite{ref:sysguide}。 & 外设硬件中断冲突、总线时钟超频，或实际带宽无法满足传感器图像吞吐。 \\
\tabbdg{9.1}& 应用与UI开发 & 基于底层 API 开发业务与图形流。需同步提供带符号表的 Debug 版供 QA 定位，及用于出货的 Release 版。 & 应用程序、源码仓库和《应用软件说明书》。 & 多线程不合理抢占导致 UI 阻塞卡顿，或发生隐蔽的长期底层内存泄漏。 \\
\tabbdg{9.2}& 算法集成 & ARM与算法组联合执行内存映射与算力加速集成。需在说明中明确算法上下文及多核/NPU资源预估。 & 由算法组同事输出算法调用 Demo 及说明。 & 算法 `.so` 库的交叉编译工具链不兼容，或实际算力消耗远超硬件芯片瓶颈。 \\
\tabbdg{10} & 板筛固件大纲 & 编写覆盖外设全要素的板级自检固件。提供给装配车间使用的操作级 SOP，降低工人的判断和操作门槛。 & 完整产品软件镜像（整包）、《测试文档》与BIT脚本等内容。 & 产测大纲存在外设硬件漏测项，导致不良品未经拦截直接流入终端客户环节。 \\
\tabbdg{11} & 黑盒集成测试 & QA独立执行边界压测、黑盒功能及 UI 健壮性验证。包含迭代优化记录与遗留已知问题（Known Issues）清单。 & 由 QA 组同事提交 Bug 反馈到相关系统。 & 异常边界条件触发极低概率的偶发 Bug，导致难以复现场景与快速出 Patch 修复。 \\
\tabbdg{12} & 环境应力(ESS) & 温箱执行 $-40^\circ\mathrm{C}\sim+85^\circ\mathrm{C}$ 极限拷机。需提供监控截图及触发死机时的 Kernel Dump 日志深度分析结论。 & 由 EE 组和 ARM 组联合出具《可靠性失效分析报告》 & 硬件晶振随温漂发生频偏导致硬死机，引发软硬件与系统各组之间的责任扯皮。 \\
\tabbdg{13} & 资料归档转产 & 工程源码/图纸合入主干受控冻结。图纸大一统，正式进入文件服务器，标志着该项目研发流程顺利闭环。 & 转产 Release 工程包、源码仓库等。 & 软硬件版本强绑定管理混乱未对齐，导致进入试产期间产线工人烧录了错误固件。 \\
\bottomrule
\end{tabularx}
\end{table}

\newpage
% =====================================================================================
% PAGE 3: 参考资料 (References)
% =====================================================================================

\renewcommand{\refname}{\large \textbf{\color{themeblue} 参考资料 (References)}}

\begin{thebibliography}{99}
    \addcontentsline{toc}{section}{参考资料}
    
    \bibitem{ref:sysguide}
    \href{run:vs8x9_sysguide.pdf}{\textcolor{blue}{\underline{\textit{VS8x9 系统指南 (Ver 1.2)}}}} [R]. 2025-10-13.
    
    \bibitem{ref:protocol}
    \href{run:protocol.pdf}{\textcolor{blue}{\underline{\textit{HTQ15-L DVP 传输协议规范 (HTQ15-L DVP Protocol Specification, Ver 1.1)}}}} [R]. 适用平台: Rockchip RK3588 / HTQ15-L. 2026-01-19.

\end{thebibliography}

\end{document}